%قبل از شروع مسئله ترجیحاً چیزی ننویسید. اما اگر خواستید بنویسید، در قالب لیستی (فقرات و شمارش) بنویسید.
%برای اضافه کردن نمره، بعدی از شروع‌مسئله، [x] را اضافه کنید و به جای x  نمره مورد نظر را قرار دهید. برای زیرمسئله‌ها هم به همین صورت است.
\شروع‌مسئله[32]
دزدی پس از دستبرد زدن به یک بانک می‌خواهد به نقطۀ امنی از شهر برای پنهان شدن برود به طوری که توسط پلیس دستگیر نشود. فرض کنید نقطۀ شروع او $s$ و نقطۀ هدف او $f$ باشد. او گراف نقاط مختلف شهر را در اختیار دارد. یال‌های این گراف نشان‌دهندۀ خیابان‌ها هستند و وزن یک یال نشان‌دهندۀ هزینۀ مورد نیاز برای طی کردن آن خیابان است که عددی مثبت است. برخی از نقاط این گراف برای دزد مخفیگاه هستند و پلیس در این مکان‌ها نمی‌تواند دزد را شناسایی کند. پلیس هر سه دقیقه یک بار به نیروهای خود در سطح شهر دستور بازرسی افراد را می‌دهد. طی کردن هر یال برای دزد یک دقیقه طول می‌کشد. دزد برای این که توسط پلیس شناسایی نشود باید هر سه دقیقه یک بار زمانی که پلیس دستور بازرسی می‌دهد در یکی از مخفیگاه‌ها باشد. دقت کنید که نقطه‌ی هدف نیز یک مخفیگاه است.

\شروع{زیرمسائل}

\زیرمسئله[12]
فرض کنید شهر به صورت یک جدول $m\times n$ باشد که خانه‌های آن رأس‌های گراف هستند و هر خانه به خانه‌های مجاور ضلعی خود یال دارد. هر خانه را با $(x, y)$ نشان می‌دهیم. نقطۀ $s$ در جنوب غربی شهر در خانۀ $(0, 0)$ قرار دارد و نقطۀ $f$ در شمال شرقی شهر در خانۀ $(m, n)$ قرار دارد. هر دقیقه دزد یا یک خانه به سمت شمال و یا یک خانه به سمت شرق می‌تواند حرکت کند. دقت کنید که دزد نمی‌تواند بعد از یک دقیقه در همان خانه باقی بماند و حتما باید حرکت کند. دزد چگونه می‌تواند مسیر با کمترین هزینه را از $s$ به $f$ پیدا کند به طوری که توسط پلیس دستگیر نشود؟ زمان الگوریتم شما باید $O(\vert V \vert +  \vert E \vert)$ باشد.

\زیرمسئله[20]
حال مسئله را در حالت کلی برای گراف دلخواه $G$ حل کنید. در این جا فرض کنید که دزد می‌تواند برای مدت دلخواهی بر حسب دقیقه در یک رأس بماند و حرکتی نکند. الگوریتمی با زمان $O(\vert V \vert \log\vert V \vert +  \vert E \vert)$ برای یافتن کم‌هزینه‌ترین مسیر از $s$ به $f$ ارائه دهید به طوری که دزد توسط پلیس دستگیر نشود.

\متن‌سیاه{راهنمایی:}
گرافی تشکیل دهید که رئوس آن موئلفه‌های دیگری علاوه بر مکان دزد داشته باشند.

\پایان{زیرمسائل}

\solution{
	
آ. از الگوریتم \متن‌لاتین{BFS} با \متن‌لاتین{Branch and Bound} استفاده می‌کنیم. به این صورت که سه لایه در \متن‌لاتین{BFS} جلو می‌رویم، با نگاه داشتن طول کوتاه‌ترین مسیر که به یک مخفیگاه مشخص در لایه‌ی سوم ختم می‌شود. این الگوریتم را تا جایی ادامه می‌دهیم که به $f$ برسیم.


ب. سه گراف $G_1$ و $G_2$ و $H$ را به این صورت می‌سازیم: برای هر رأس $v$ در گراف اصلی یک رأس $v_1$ در $G_1$ و یک رأس $v_2$ در $G_2$ در نظر می‌گیریم. همچنین برای هر مخفیگاه $v$ در گراف اصلی یک رأس $v_h$ در $H$ در نظر می‌گیریم. برای هر یال $u \ v$ در گراف اصلی یک یال از $u_1$ به $v_2$ در نظر می‌گیریم. در صورتی که $v$ یک مخفیگاه باشد یک یال از $u_2$ به $v_h$ در نظر می‌گیریم. در صورتی که $u$ یک مخفیگاه باشد یک یال از $u_h$ به $v_1$ در نظر می‌گیریم. یک رأس جدید به نام $F$ در نظر می‌گیریم و رئوس $f_1$ و $f_2$ و $f_h$ را با یال به طول $0$ به آن وصل می‌کنیم.

حال کافی است کوتاه‌ترین مسیر را از $s_1$ به $F$ پیدا کنیم.

}




%%%%%%%%%%%%%%%%%%%%%%%%%%%%%%%%%%%%%%%%%%%%%%%%%%%%%
% برخی دستورات مفید برای تمرین‌ها. %
%%%%%%%%%%%%%%%%%%%%%%%%%%%%%%%%%%%%%%%%%%%%%%%%%%%%%

% فرمول‌های نمونه:
%$f(x) = \Theta(x)$
%$T(x, y) \leq \log(x) + 2^y + \binom{2n}{n}$
% {\tt code\_function}


% لیست موارد بدون شماره‌گذاری:
%\شروع{فقرات}
%    \فقره مورد اول 
%        \شروع{فقرات}
%            \فقره مورد اول داخلی 
%                \شروع{فقرات}
%                    \فقره اولین مورد خیلی داخلی 
%                    \فقره دومین مورد خیلی داخلی 
%                \پایان{فقرات}
%            \فقره مورد دوم داخلی 
%        \پایان{فقرات}
%    \فقره مورد دوم 
%\پایان{فقرات}

% لیست موارد شماره‌گذاری شده:
%\شروع{شمارش}
%    \فقره یک مورد 
%    \فقره یک مورد دیگر 
%    \فقره باز هم یک مورد دیگر
%\پایان{شمارش}

% برای نوشتن یک فرمول چند خطی که یک شماره بخوره. برای حالت‌های دیگر به سایت زیر مراجعه کنید:
% https://www.overleaf.com/learn/latex/Aligning_equations_with_amsmath
%\begin{align} 
%    \begin{split}
%        (x+y)^3 &= (x+y)^2(x+y) \\
%                &= (x^2+2xy+y^2)(x+y) \\
%                &= (x^3+2x^2y+xy^2) + (x^2y+2xy^2+y^3) \\
%                &= x^3+3x^2y+3xy^2+y^3
%    \end{split}                                 
%\end{align}

% نوشتن ماتریس و آرایه دوبعدی
%\begin{align}
%    A = 
%    \begin{bmatrix}
%        A_{11} & A_{21} \\
%        A_{21} & A_{22}
%    \end{bmatrix}
%\end{align}

% اضافه کردن تصویر با شرح. برای اضافه کردن تصویر بدون شرح، فقط دستور درج‌تصویر را استفاده کنید.
%\شروع{شکل}[h!]
%\تنظیم‌ازوسط
%  \درج‌تصویر[پهنا=0.5\پهنای‌متن]{img.jpg}
%\شرح{شرح تصویر}
%\برچسب{تصویر:۱}
%\پایان{شکل}
%\پایان{زیرمسائل}

%کد زیر را ببینید.
%%\lstinputlisting{source_filename.py}
%\begin{latin}
%\begin{lstlisting}
%# Import the modules
%import sys
%import random
%
%ans = True
%
%while ans:
%question = raw_input("Ask the magic 8 ball a question: (press enter to quit) ")
%
%answers = random.randint(1,8)
%
%if question == "":
%sys.exit()
%
%elif answers == 1:
%print "It is certain"
%
%elif answers == 2:
%print "Outlook good"
%
%elif answers == 3:
%print "You may rely on it"
%
%elif answers == 4:
%print "Ask again later"
%
%elif answers == 5:
%print "Concentrate and ask again"
%
%elif answers == 6:
%print "Reply hazy, try again"
%
%elif answers == 7:
%print "My reply is no"
%
%elif answers == 8:
%print "My sources say no"
%\end{lstlisting}
%\end{latin}
